\chapter{Special Elements}
%
%
The Schema definition of the XML parameter file currently allows for two
special elements to be specified. These are
%
%
\begin{itemize}
\item \xml{amExpert}
\item \xml{devel}
\end{itemize}
%
%
These are first level elements coming at the end of the parameter file, see
Sect.~\ref{sec:general}.

The \xml{amExpert} is just an empty element without attributes. Specifying
\xml{amExpert} in the parameter file turns of the Expert() module of the
CFS++ --- OLAS interface. This expert module normally tries to set certain
parameters to sensible values depending on the values of other parameters or
to fix inconsistencies in the input parameters. E.g.~the expert module will
turn off preconditioning, if a direct solver is specified, or change the
matrix type to \xml{hypreMatrix}, if one of the solvers of the hypre library
is specified. There may be situations, however, where this behaviour is not
desired, e.g.~one may want to test the performance of an ILU preconditioner
without the \xml{Sloan} matrix re-ordering the expert would trigger by default.
Note that turning off the expert module by specifying \xml{amExpert} may lead
to run-time aborts in OLAS due to inconsistent settings. So be sure to really
be an expert, if you do this.

The \xml{devel} element can have arbitrary content, but no children or
attributes. It is solely there for the sake of the developer who wants to
wants to quickly test something, without the need to alter the Schema
definition itself. For normal users the element is not useful, since it is
not read in the standard CFS++ source code.
